\section{Contoh Terhitung}

\begin{example}[Algoritma Euclidean dan invers modulo]
\begin{enumerate}
    \item \textbf{PBB dan koefisien B\'ezout dari pasangan $(1970,1066)$.} Terapkan algoritma pembagian berulang berikut untuk memperoleh sisa terakhirnya:
    \[
    \begin{aligned}
    1970 &= 1 \cdot 1066 + 904\\
    1066 &= 1 \cdot 904 + 162\\
    904  &= 5 \cdot 162 + 94\\
    162  &= 1 \cdot 94 + 68\\
    94   &= 1 \cdot 68 + 26\\
    68   &= 2 \cdot 26 + 16\\
    26   &= 1 \cdot 16 + 10\\
    16   &= 1 \cdot 10 + 6\\
    10   &= 1 \cdot 6 + 4\\
    6    &= 1 \cdot 4 + 2\\
    4    &= 2 \cdot 2 + 0
    \end{aligned}
    \]
    Sisa bukan nol terakhir adalah $2$, jadi $\text{PBB}(1970,1066) = 2$. Substitusi balik memberi koefisien B\'ezout
    \[ 1970 \cdot (-204) + 1066 \cdot 377 = 2. \]
    \item \textbf{Invers modulo $550 \bmod 1769$.} Karena $\text{PBB}(550,1769) = 1$, inversnya ada. Uji nilai $x = 550$:
    \[ 550 \cdot 550 = 302\,500 = 171 \cdot 1769 + 1, \]
    sehingga $550 \cdot 550 \equiv 1 \pmod{1769}$, yaitu $550^{-1} \equiv 550 \pmod{1769}$.
\end{enumerate}
\end{example}

\begin{example}[Entropi Shannon sebuah distribusi]
Pesan $X = \code{AABBCBDB}$ terdiri atas 8 simbol dengan peluang $p(A)=2/8$, $p(B)=4/8$, $p(C)=1/8$, dan $p(D)=1/8$. Substitusikan ke rumus entropi:
\[
\begin{aligned}
H(X) &= -\sum_{i} p(x_i) \log_2 p(x_i)\\
     &= -\left(\tfrac14 \log_2 \tfrac14 + \tfrac12 \log_2 \tfrac12
        + \tfrac18 \log_2 \tfrac18 + \tfrac18 \log_2 \tfrac18\right)\\
     &= -\left(\tfrac14(-2) + \tfrac12(-1) + \tfrac18(-3) + \tfrac18(-3)\right)\\
     &= 1{,}75\ \text{bit/simbol}.
\end{aligned}
\]
Karena entropi maksimum untuk 4 simbol adalah $\log_2 4 = 2$ bit/simbol, distribusi yang tidak seragam ini menghasilkan entropi yang lebih rendah dan karena itu lebih mudah dianalisis secara statistik.
\end{example}

\begin{example}[Fungsi totient Euler dan Teorema Euler]
Hitung $\varphi(20)$, yaitu banyaknya bilangan pada rentang $1 \le k \le 20$ yang
relatif prima terhadap 20. Bilangan-bilangan itu adalah
\[ 1,\ 3,\ 7,\ 9,\ 11,\ 13,\ 17,\ 19, \]
seluruhnya delapan buah, jadi $\varphi(20) = 8$. Pemeriksaan melalui sifat
multiplikatif memberi hasil yang sama: karena $20 = 4 \cdot 5$ dan
$\text{PBB}(4,5)=1$,
\[ \varphi(20) = \varphi(4) \cdot \varphi(5) = 2 \cdot 4 = 8. \]

Selanjutnya periksa Teorema Euler untuk $a = 3$ dan $n = 10$. Karena
$\text{PBB}(3,10) = 1$ dan $\varphi(10) = \varphi(2)\cdot\varphi(5) = 1 \cdot 4 = 4$,
teorema menyatakan $3^4 \equiv 1 \pmod{10}$. Pemeriksaannya:
\[ 3^4 = 81 = 8 \cdot 10 + 1, \]
sehingga benar bahwa $3^4 \equiv 1 \pmod{10}$. Perhatikan bahwa yang diperiksa bukan
nilai $3^4$ itu sendiri, melainkan sisanya ketika dibagi 10.
\end{example}

\begin{example}[Akar primitif dan logaritma diskrit]
Tinjau modulo $p = 7$. Untuk $g = 3$, hitung pangkat-pangkatnya secara berurutan:
\[
\begin{aligned}
3^1 &= 3 \equiv 3, &\qquad 3^2 &= 9 \equiv 2,\\
3^3 &= 27 \equiv 6, &\qquad 3^4 &= 81 \equiv 4,\\
3^5 &= 243 \equiv 5, &\qquad 3^6 &= 729 \equiv 1.
\end{aligned}
\]
Siklusnya menghasilkan $3, 2, 6, 4, 5, 1$ --- keenam anggota $\{1,2,3,4,5,6\}$ muncul
tepat satu kali, sehingga $g = 3$ adalah akar primitif modulo 7.

Sebagai pembanding, untuk $g = 2$:
\[ 2^1 \equiv 2, \qquad 2^2 \equiv 4, \qquad 2^3 = 8 \equiv 1, \]
sehingga siklusnya hanya berpanjang tiga dan hanya menghasilkan $\{1,2,4\}$. Karena
tidak seluruh anggota $\{1,\dots,6\}$ tercakup, $g = 2$ bukan akar primitif modulo 7.

Selanjutnya tinjau logaritma diskrit $7^x \equiv 15 \pmod{41}$. Hitung pangkat 7
secara berurutan sambil mengurangi modulo 41:
\[
\begin{aligned}
7^1 &= 7 \equiv 7, &\qquad 7^2 &= 49 = 41 + 8 \equiv 8,\\
7^3 &= 7 \cdot 8 = 56 = 41 + 15 \equiv 15.
\end{aligned}
\]
Jadi $x = 3$. Perhatikan bahwa menghitung $7^3 \bmod 41$ hanya menuntut dua perkalian,
sedangkan menemukan $x = 3$ dari $7^x \equiv 15$ pada modulus yang berukuran realistis
--- misalnya 2048 bit --- tidak memiliki cara penyelesaian yang praktis.
\end{example}

\begin{example}[Entropi plainteks berita Kota Ternate]
Plainteks berita Kota Ternate pada Bagian~\ref{sec:teori-informasi} memuat 296 huruf
setelah spasi dan tanda baca dibuang, dengan 21 huruf berbeda. Enam huruf yang paling
sering muncul adalah:

\begin{center}
\begin{tabular}{@{}lcccccc@{}}
\hline
Huruf & A & N & E & I & K & T \\
Cacah & 59 & 34 & 29 & 22 & 17 & 17 \\
Peluang & $0{,}19932$ & $0{,}11486$ & $0{,}09797$ & $0{,}07432$ & $0{,}05743$ & $0{,}05743$ \\
\hline
\end{tabular}
\end{center}

Kontribusi keenam huruf ini terhadap entropi dihitung dengan suku $-p \log_2 p$:
\[
\begin{aligned}
-p(A)\log_2 p(A) &= 0{,}19932 \times 2{,}327 &= 0{,}4638\\
-p(N)\log_2 p(N) &= 0{,}11486 \times 3{,}122 &= 0{,}3586\\
-p(E)\log_2 p(E) &= 0{,}09797 \times 3{,}352 &= 0{,}3284\\
-p(I)\log_2 p(I) &= 0{,}07432 \times 3{,}750 &= 0{,}2787\\
-p(K)\log_2 p(K) &= 0{,}05743 \times 4{,}122 &= 0{,}2367\\
-p(T)\log_2 p(T) &= 0{,}05743 \times 4{,}122 &= 0{,}2367
\end{aligned}
\]
Jumlah keenamnya adalah $1{,}9029$ bit. Lima belas huruf sisanya --- yang masing-masing
muncul jauh lebih jarang --- menyumbang $1{,}9959$ bit lagi, sehingga totalnya
\[ H(P) = 1{,}9029 + 1{,}9959 = 3{,}8988\ \text{bit/karakter}. \]
Dibandingkan dengan entropi maksimum alfabet 26 huruf, pencapaiannya
\[ \frac{3{,}8988}{4{,}7004} \times 100\% \approx 82{,}95\%. \]
Perhatikan bahwa lima belas huruf yang jarang muncul justru menyumbang bagian entropi
yang lebih besar daripada enam huruf tersering: inilah wujud nyata dari tafsir entropi
sebagai ``rata-rata kejutan''. Huruf yang jarang mengagetkan, dan kejutan itulah yang
mengandung informasi.
\end{example}
